% !TeX program = xelatex
% ============================================================================
%  第 9 课　D₄ 的子群与二面体群的关系　—— 教师版讲义
%  与 02-课件/第09课-D4的子群/第09课-D4的子群.tex 同步
% ============================================================================

\jhlesson{9}{\texorpdfstring{\(D_4\)}{D4} 的子群与二面体群的关系}{时长 45 分钟\quad·\quad 教具：正方形纸片\quad·\quad 本课不发单页}

\begin{mubiao}
  \begin{itemize}
    \item 能找出 \(D_4\) 里面能自成一体的小群（子群）。
    \item 能说出旋转与翻折之间那条特殊关系：\(fr=r^{-1}f\)。
    \item 知道 \(D_n\) 这个名字叫\textbf{二面体群}，知道正 \(n\) 边形有 \(2n\) 个对称动作。
  \end{itemize}
\end{mubiao}

\noindent{\small 本课节奏：0—6 回顾｜6—26 找 \(D_4\) 的子群｜26—40 旋转与翻折的关系｜40—45 收尾。全程至少 4 次「用手指回答」。}

\section*{一、回顾（0—6 分钟）}

把上节课在 \(S_3\) 里找子群的办法再写一遍：圈出一小堆动作，把它们的复合填成一张表；
结果全在里面、没有跑出去，它就是子群。

一句话过渡：今天把这条办法照搬到正方形上。

\section*{二、找 \(D_4\) 的子群（6—26 分钟）}

\begin{caozuo}
  一步一步带着全班圈，每圈一个就填一张小表：
  \begin{enumerate}
    \item \(H=\{e,\ r^2\}\)：\(r^2\) 的阶是 \(2\)，表里没跑出去，\textbf{是}子群，大小 \(2\)；
    \item \(H=\{e,\ r,\ r^2,\ r^3\}\)：四个转动转来转去还是转，\textbf{是}子群，大小 \(4\)；
    \item \(H=\{e,\ f\}\)：翻两次回来，\textbf{是}子群；同样的圈子有 \(4\) 个
      （\(f,\ rf,\ r^2f,\ r^3f\) 各一个），每个大小 \(2\)；
    \item \(H=\{e,\ r^2,\ f,\ r^2f\}\)（含两条对角线）：填表后全在里面，\textbf{是}子群，大小 \(4\)；
    \item \(H=\{e,\ r^2,\ rf,\ r^3f\}\)（含两条中线）：也一样，\textbf{是}子群，大小 \(4\)。
  \end{enumerate}
  再让全班自己试几个不成的，比如 \(\{e,r\}\)：\(r\cdot r=r^2\) 不在里面，跑出去了。
\end{caozuo}

把结果汇总成一张表：

\begin{center}
\begin{tabular}{@{}clc@{}}
  \toprule
  大小 & 子群 & 个数 \\
  \midrule
  \(1\) & \(\{e\}\) & \(1\) \\
  \(2\) & \(\{e,r^2\}\)；\(\{e,f\}\)；\(\{e,rf\}\)；\(\{e,r^2f\}\)；\(\{e,r^3f\}\) & \(5\) \\
  \(4\) & \(\{e,r,r^2,r^3\}\)；\(\{e,r^2,f,r^2f\}\)；\(\{e,r^2,rf,r^3f\}\) & \(3\) \\
  \(8\) & \(D_4\)（全部八个） & \(1\) \\
  \bottomrule
\end{tabular}
\end{center}

一共 \(1+5+3+1=10\) 个子群。大小 \(1,2,4,8\) 都能整除 \(8\)。

\begin{zhu}
  观察就好，不要证明。第 10 课会在模 \(12\) 的加法里再看一次同样的现象，
  那时候学生自己就能说出「小圈的大小整除大圈的大小」。
\end{zhu}

\begin{caozuo}
  手指回答：\(\{e,r^2\}\) 是子群吗？（是）
  \(\{e,r\}\) 是子群吗？（不是，\(r\cdot r=r^2\) 跑出去了）
  一个大小是 \(8\) 的群，子群大小可能是 \(3\) 吗？（不可能，\(3\) 不能整除 \(8\)）
\end{caozuo}

\section*{三、旋转与翻折的关系（26—40 分钟）}

\begin{caozuo}
  拿纸片走两条路，两个终点用不同颜色标出来：
  \begin{itemize}
    \item 左边：先翻（\(f\)），再转 \(90^\circ\)（\(r\)）——走的是 \(fr\)；
    \item 右边：先转 \(270^\circ\)（\(r^3\)），再翻（\(f\)）——走的是 \(r^3f\)。
  \end{itemize}
  两边的终点落在同一个位置，于是板书：
  \[
    \mathbf{fr=r^{-1}f} .
  \]
\end{caozuo}

白话念一遍：「先翻后转，等于先\textbf{倒着}转、再翻。」
在正方形里 \(r\) 的倒着做是 \(r^3\)，所以这个式子也写成 \(fr=r^3f\)。

\begin{dingli}[二面体群]
  正 \(n\) 边形的全部对称动作，按复合构成一个群，叫做\textbf{二面体群}，记作 \(D_n\)，
  一共有 \(2n\) 个动作：\(n\) 个转动加 \(n\) 个翻折。
\end{dingli}

\begin{zhu}
  三角形的 \(D_3\) 就是我们前面写的 \(S_3\)（都是 \(6\) 个动作）。
  这里不必展开为什么两个记号都有人用，点到为止。
\end{zhu}

\section*{四、收尾（40—45 分钟）}

板书两句话：\(D_4\) 里能找到 \(10\) 个子群，大小都能整除 \(8\)；
旋转与翻折之间有一条关系 \(fr=r^{-1}f\)。

收尾伏笔：\(8\) 个动作，其实只要 \(r\) 和 \(f\) 两个就够生成了。

\begin{xiaojie}
  \begin{itemize}
    \item \(D_4\) 有 \(10\) 个子群，大小是 \(1,2,4,8\)，都能整除 \(8\)。
    \item 关系式 \(fr=r^{-1}f\)，也就是 \(fr=r^3f\)。
    \item \(D_n\) 叫二面体群，正 \(n\) 边形共有 \(2n\) 个对称动作。
  \end{itemize}
\end{xiaojie}

\noindent\textbf{下节课}：生成元、循环群与循环群的子群——\(8\) 个动作为什么两个就够。